%%%%%%%%%%%%%%%%%%%%%%%%%%%%%%%%%
%
%       EXERCÍCIO
%
%%%%%%%%%%%%%%%%%%%%%%%%%%%%%%%%%

\definevalorquestao

\begin{exercicioBanco}[\valorquestao]
Para cada um dos casos abaixo, escreva o espaço amostral correspondente e conte seus elementos.
%
\begin{enumerate}[(a)]
    \item Uma moeda é lançada duas vezes e observam-se as faces obtidas.
    %
    \item Um dado é lançado duas vezes e a ocorrência de face par ou ímpar é observada.
    %
    \item Uma urna contém 10 bolas azuis e 10 bolas vermelhas com dimensões iguais. Três bolas são selecionas ao acaso com reposição e as cores são anotadas.
    %
    \item Dois dados são lançados simultaneamente e estamos interessados na soma das faces observadas.
\end{enumerate}
%
\end{exercicioBanco}

